% Slide 5: Anatomy of a Great Spec
% Visual: A mock Markdown spec document showing the key sections, with annotations
%   highlighting what each section gives the AI.
% Talking Points:
%   - A spec has 6 essential sections: Overview, Data Model, Functional Requirements,
%     API Contracts, Non-Functional Requirements, and Acceptance Criteria.
%   - Each section constrains the AI in a specific, useful way.
%   - The spec shown here is the one we will feed to Claude Code in the live demo.
% Presenter Script:
%   "Let's look at what a real spec looks like. This is the actual document
%    I'm about to feed to Claude Code in the demo. It has six sections. The
%    Overview tells the AI WHAT we're building and WHY -- this prevents scope
%    creep. The Data Model defines every entity and relationship -- no guessing.
%    Functional Requirements are numbered, testable statements: FR-001, the
%    system SHALL do X. API Contracts define every endpoint, request, response,
%    and error code. Non-functional requirements cover performance, security,
%    accessibility. And Acceptance Criteria are the definition of done -- the AI
%    can self-verify against these. Notice: this is Markdown. Not UML, not a
%    proprietary format. It lives in your repo, gets versioned with git, and
%    every engineer can read it."

\section{The Spec}

\begin{frame}[fragile]{Anatomy of a Great Spec}
  \vspace{-0.6cm}
  \begin{columns}[T]
    \begin{column}{0.52\textwidth}
      \begin{lstlisting}[style=markdown, basicstyle=\ttfamily\fontsize{4}{5}\selectfont]
# Alert Dashboard - SRS

## 1. Overview
Real-time alert dashboard for
monitoring Keylime attestation events.

## 2. Data Model
- Alert: {id, severity, agent_id,
  message, timestamp, acknowledged}
- Severity: enum(critical, warning, info)

## 3. Functional Requirements
- FR-001: Display alerts in real-time
- FR-002: Filter by severity and date
- FR-003: Acknowledge/dismiss alerts
- FR-004: Paginate results (25/page)

## 4. API Contracts
GET /api/alerts?severity=critical
Response: {alerts: [], total, page}
Errors: 400, 401, 500

## 5. Non-Functional Requirements
- NFR-001: < 200ms response time
- NFR-002: WCAG 2.1 AA accessible
- NFR-003: 95% test coverage

## 6. Acceptance Criteria
- [ ] Alert list renders in < 1s
- [ ] Filters reduce result set
- [ ] Ack persists across reload
      \end{lstlisting}
    \end{column}
    \begin{column}{0.48\textwidth}
      \vspace{0.2cm}
      \begin{tikzpicture}[scale=0.65, transform shape,
          annot/.style={rectangle, draw, rounded corners, fill=keylimeblue!8,
                        text width=4.6cm, font=\scriptsize, inner sep=0.15cm}]

        \node[annot] (a1) at (0,4.4) {
          \faIcon{compass}~\textbf{Overview}\\
          Prevents scope creep -- tells AI the boundaries of the world
        };
        \node[annot] (a2) at (0,2.9) {
          \faIcon{database}~\textbf{Data Model}\\
          Eliminates guessing about entity shapes and relationships
        };
        \node[annot] (a3) at (0,1.4) {
          \faIcon{clipboard-list}~\textbf{Functional Reqs}\\
          Numbered, testable: \texttt{FR-NNN}\\
          Each one maps to test cases
        };
        \node[annot] (a4) at (0,-0.1) {
          \faIcon{plug}~\textbf{API Contracts}\\
          Endpoints, payloads, error codes -- no hallucinated APIs
        };
        \node[annot] (a5) at (0,-1.6) {
          \faIcon{tachometer-alt}~\textbf{Non-Functional Reqs}\\
          Performance, security, a11y -- constraints, not wishes
        };
        \node[annot] (a6) at (0,-3.1) {
          \faIcon{check-double}~\textbf{Acceptance Criteria}\\
          Definition of ``done'' -- AI self-verifies against these
        };

      \end{tikzpicture}
    \end{column}
  \end{columns}

  \vspace{-0.3cm}

  \begin{block}{\faIcon{lightbulb}~NOTE for the Product Owner}
    \centering\scriptsize
    Could this be mapped to JIRA Epics?
  \end{block}
\end{frame}
